% GENERATED FILE. Edit canonical lessons, then run npm run generate:books.
% canonical-source-hash: fnv1a64:e9d9d920c6922d59
% canonical-lessons: RU-C344-logika, RU-C344-logichnyy, RU-C344-vyvod, RU-C344-rassuzhdat, RU-C344-ochevidnyy, RU-R344-review-337-344

\chapter{Logic and conclusions}
\label{ch:ru-logic-and-conclusions}
\hlchaptermodality{344}

\begin{chapteropening}
\textbf{By the end of this chapter:} \emph{I can follow the logic of an argument, say whether it is logical, draw a conclusion, reason aloud and call something obvious.}

Complete the last lesson of chapter 344: I can follow the logic of an argument, say whether it is logical, draw a conclusion, reason aloud and call something obvious.
\end{chapteropening}

\section[\texorpdfstring{\ru{логика} (lógika)}{логика (lógika)}]{\ru{логика} (lógika) — logic}

\label{lesson:RU-C344-logika}



\begin{quote}
Before the new one: say the Russian for to condemn, then the Russian for to reproach.
\end{quote}

\subsection*{You'll want to know: \ru{логика}}
\textbf{\ru{логика}} (\emph{lógika}) — \textquotedblleft{}logic\textquotedblright{}.

\begin{quote}
\ru{В} \ru{его} \ru{словах} \ru{нет} \ru{логики}.
\end{quote}

\begin{quote}
\emph{V yegó slovákh net lógiki.}
\end{quote}

\begin{quote}
There is no logic in what he says.
\end{quote}

Feminine. \emph{\ru{По} \ru{логике}} — logically speaking.

\begin{cousinweb}
\textbf{\ru{логика}} is usually traced, through western languages, to Greek \emph{logik\'{\={e}}} \textquotedblleft{}the art of reasoning\textquotedblright{}, built on \emph{lógos} \textquotedblleft{}word, reason\textquotedblright{}.
\end{cousinweb}

\begin{grammarlens}[title={What you've built}]
One more word for telling, explaining and describing at B1.
\end{grammarlens}

\subsection*{Your turn}
\begin{itemize}
\raggedright
  \item \emph{Say it:} \emph{lógika}
  \item \emph{Say it:} \emph{lógika}, once more
  \item \emph{Say it:} a sentence of your own with \emph{lógika}
  \item \emph{Recall:} say the Russian for to condemn, then the Russian for to reproach, then say \emph{lógika} again
\end{itemize}

\begin{tcolorbox}[breakable,colback=teal!4,colframe=teal!35!black,title={Before you move on}]
What does \ru{логика} mean? (\textquotedblleft{}Logic\textquotedblright{}.) Say it once more.
\end{tcolorbox}
\section[\texorpdfstring{\ru{логичный} (logíchnyy)}{логичный (logíchnyy)}]{\ru{логичный} (logíchnyy) — logical}

\label{lesson:RU-C344-logichnyy}



\begin{quote}
Before the new one: say the Russian for to reproach, then the Russian for logic.
\end{quote}

\subsection*{You'll want to know: \ru{логичный}}
\textbf{\ru{логичный}} (\emph{logíchnyy}) — \textquotedblleft{}logical\textquotedblright{}.

\begin{quote}
\ru{Твой} \ru{план} \ru{вполне} \ru{логичен}.
\end{quote}

\begin{quote}
\emph{Tvoy plan vpolné logíchen.}
\end{quote}

\begin{quote}
Your plan is quite logical.
\end{quote}

Short forms \emph{\ru{логичен}, \ru{логична}}. Of people too: \emph{\ru{логичный} \ru{человек}}.

\begin{cousinweb}
\textbf{\ru{логичный}} is built on \emph{\ru{логика}} with the Russian adjective suffix \emph{-\ru{ичный}}; \emph{\ru{логика}} is usually traced to Greek \emph{logik\'{\={e}}}, on \emph{lógos} \textquotedblleft{}word, reason\textquotedblright{}.
\end{cousinweb}

\begin{grammarlens}[title={What you've built}]
One more word for telling, explaining and describing at B1.
\end{grammarlens}

\subsection*{Your turn}
\begin{itemize}
\raggedright
  \item \emph{Say it:} \emph{logíchnyy}
  \item \emph{Say it:} \emph{logíchnyy}, once more
  \item \emph{Say it:} a sentence of your own with \emph{logíchnyy}
  \item \emph{Recall:} say the Russian for to reproach, then the Russian for logic, then say \emph{logíchnyy} again
\end{itemize}

\begin{tcolorbox}[breakable,colback=teal!4,colframe=teal!35!black,title={Before you move on}]
What does \ru{логичный} mean? (\textquotedblleft{}Logical\textquotedblright{}.) Say it once more.
\end{tcolorbox}
\section[\texorpdfstring{\ru{вывод} (vývod)}{вывод (vývod)}]{\ru{вывод} (vývod) — a conclusion}

\label{lesson:RU-C344-vyvod}



\begin{quote}
Before the new one: say the Russian for logic, then the Russian for logical.
\end{quote}

\subsection*{You'll want to know: \ru{вывод}}
\textbf{\ru{вывод}} (\emph{vývod}) — \textquotedblleft{}a conclusion\textquotedblright{}.

\begin{quote}
\ru{Какой} \ru{вывод} \ru{можно} \ru{сделать}?
\end{quote}

\begin{quote}
\emph{Kakóy vývod mózhno sdélat?}
\end{quote}

\begin{quote}
What conclusion can we draw?
\end{quote}

Masculine. \emph{\ru{Сделать} \ru{вывод}} — to draw a conclusion. It also means withdrawal: \emph{\ru{вывод} \ru{войск}}.

\begin{cousinweb}
\textbf{\ru{вывод}} is built from \emph{\ru{вы}-} \textquotedblleft{}out\textquotedblright{} and \emph{\ru{вод}-}, the root of \emph{\ru{водить}} \textquotedblleft{}to lead\textquotedblright{}, from Proto-Slavic \emph{*voditi} \textquotedblleft{}to lead, to conduct\textquotedblright{} (EtymDB). A conclusion is what you lead out of the facts.
\end{cousinweb}

\begin{grammarlens}[title={What you've built}]
One more word for telling, explaining and describing at B1.
\end{grammarlens}

\subsection*{Your turn}
\begin{itemize}
\raggedright
  \item \emph{Say it:} \emph{vývod}
  \item \emph{Say it:} \emph{vývod}, once more
  \item \emph{Say it:} a sentence of your own with \emph{vývod}
  \item \emph{Recall:} say the Russian for logic, then the Russian for logical, then say \emph{vývod} again
\end{itemize}

\begin{tcolorbox}[breakable,colback=teal!4,colframe=teal!35!black,title={Before you move on}]
What does \ru{вывод} mean? (\textquotedblleft{}A conclusion\textquotedblright{}.) Say it once more.
\end{tcolorbox}
\section[\texorpdfstring{\ru{рассуждать} (rassuzhdát)}{рассуждать (rassuzhdát)}]{\ru{рассуждать} (rassuzhdát) — to reason, to think aloud}

\label{lesson:RU-C344-rassuzhdat}



\begin{quote}
Before the new one: say the Russian for logical, then the Russian for a conclusion.
\end{quote}

\subsection*{You'll want to know: \ru{рассуждать}}
\textbf{\ru{рассуждать}} (\emph{rassuzhdát}) — \textquotedblleft{}to reason, to think aloud\textquotedblright{}.

\begin{quote}
\ru{Он} \ru{любит} \ru{рассуждать} \ru{о} \ru{жизни}.
\end{quote}

\begin{quote}
\emph{On lyúbit rassuzhdát o zhízni.}
\end{quote}

\begin{quote}
He likes to reason about life.
\end{quote}

Imperfective only. \emph{\ru{Рассуждать} \ru{о}} + prepositional.

\begin{cousinweb}
\textbf{\ru{рассуждать}} is built from \emph{\ru{рас}-} \textquotedblleft{}apart\textquotedblright{} and \emph{\ru{судить}} \textquotedblleft{}to judge\textquotedblright{}, on \emph{\ru{суд}}, inherited from Proto-Slavic \emph{*sǫdŭ} (EtymDB). To reason is to take a question apart and judge it.
\end{cousinweb}

\begin{grammarlens}[title={What you've built}]
One more word for telling, explaining and describing at B1.
\end{grammarlens}

\subsection*{Your turn}
\begin{itemize}
\raggedright
  \item \emph{Say it:} \emph{rassuzhdát}
  \item \emph{Say it:} \emph{rassuzhdát}, once more
  \item \emph{Say it:} a sentence of your own with \emph{rassuzhdát}
  \item \emph{Recall:} say the Russian for logical, then the Russian for a conclusion, then say \emph{rassuzhdát} again
\end{itemize}

\begin{tcolorbox}[breakable,colback=teal!4,colframe=teal!35!black,title={Before you move on}]
What does \ru{рассуждать} mean? (\textquotedblleft{}To reason, to think aloud\textquotedblright{}.) Say it once more.
\end{tcolorbox}
\section[\texorpdfstring{\ru{очевидный} (ochevídnyy)}{очевидный (ochevídnyy)}]{\ru{очевидный} (ochevídnyy) — obvious}

\label{lesson:RU-C344-ochevidnyy}



\begin{quote}
Before the new one: say the Russian for a conclusion, then the Russian for to reason.
\end{quote}

\subsection*{You'll want to know: \ru{очевидный}}
\textbf{\ru{очевидный}} (\emph{ochevídnyy}) — \textquotedblleft{}obvious\textquotedblright{}.

\begin{quote}
\ru{Ответ} \ru{был} \ru{очевиден} \ru{для} \ru{всех}.
\end{quote}

\begin{quote}
\emph{Otvét byl ochevíden dlya vsekh.}
\end{quote}

\begin{quote}
The answer was obvious to everyone.
\end{quote}

Short forms \emph{\ru{очевиден}, \ru{очевидна}}. \emph{\ru{Очевидно}} — evidently, clearly.

\begin{cousinweb}
\textbf{\ru{очевидный}} is built from \emph{\ru{очи}}, the old plural of \emph{\ru{око}} \textquotedblleft{}eye\textquotedblright{}, and \emph{\ru{вид}-}, the root of \emph{\ru{видеть}} \textquotedblleft{}to see\textquotedblright{}, from Proto-Slavic \emph{*viděti} (EtymDB). What is obvious is seen with your own eyes.
\end{cousinweb}

\begin{grammarlens}[title={What you've built}]
One more word for telling, explaining and describing at B1.
\end{grammarlens}

\subsection*{Your turn}
\begin{itemize}
\raggedright
  \item \emph{Say it:} \emph{ochevídnyy}
  \item \emph{Say it:} \emph{ochevídnyy}, once more
  \item \emph{Say it:} a sentence of your own with \emph{ochevídnyy}
  \item \emph{Recall:} say the Russian for a conclusion, then the Russian for to reason, then say \emph{ochevídnyy} again
\end{itemize}

\begin{tcolorbox}[breakable,colback=teal!4,colframe=teal!35!black,title={Before you move on}]
What does \ru{очевидный} mean? (\textquotedblleft{}Obvious\textquotedblright{}.) Say it once more.
\end{tcolorbox}
\section[(dialogue)]{Reasons and opinions — a review}

\label{lesson:RU-R344-review-337-344}



\begin{quote}
Say the Russian for a point of view and a position.
\end{quote}

\subsection*{Your turn}
Say these aloud:

\begin{itemize}
\raggedright
  \item a point of view, a position, a view, to declare, to be considered
  \item an argument, a point made, to prove, proof, convincing
  \item to convince, to change someone’s mind, a conviction, confidence, to doubt
  \item to agree, agreement, disagreement, to object, an objection
  \item to contradict, to reject, a compromise, to insist, to admit
  \item to discuss, a discussion, a debate, to speak one’s mind, to share
  \item to criticise, criticism, to approve of, to condemn, to reproach
  \item logic, logical, a conclusion, to reason, obvious
\end{itemize}

\begin{tcolorbox}[breakable,colback=teal!4,colframe=teal!35!black,title={Before you move on}]
A point of view? (\textbf{\ru{точка} \ru{зрения}}.) A position? (\textbf{\ru{позиция}}.) A view? (\textbf{\ru{взгляд}}.) To declare? (\textbf{\ru{заявлять}}.) To be considered? (\textbf{\ru{считаться}}.) An argument? (\textbf{\ru{аргумент}}.) A point made? (\textbf{\ru{довод}}.) To prove? (\textbf{\ru{доказывать}}.) Proof? (\textbf{\ru{доказательство}}.) Convincing? (\textbf{\ru{убедительный}}.) To convince? (\textbf{\ru{убеждать}}.) To change someone’s mind? (\textbf{\ru{переубеждать}}.) A conviction? (\textbf{\ru{убеждение}}.) Confidence? (\textbf{\ru{уверенность}}.) To doubt? (\textbf{\ru{сомневаться}}.) To agree? (\textbf{\ru{соглашаться}}.) Agreement? (\textbf{\ru{согласие}}.) Disagreement? (\textbf{\ru{несогласие}}.) To object? (\textbf{\ru{возражать}}.) An objection? (\textbf{\ru{возражение}}.) To contradict? (\textbf{\ru{противоречить}}.) To reject? (\textbf{\ru{отвергать}}.) A compromise? (\textbf{\ru{компромисс}}.) To insist? (\textbf{\ru{настаивать}}.) To admit? (\textbf{\ru{признавать}}.) To discuss? (\textbf{\ru{обсуждать}}.) A discussion? (\textbf{\ru{обсуждение}}.) A debate? (\textbf{\ru{дискуссия}}.) To speak one’s mind? (\textbf{\ru{высказываться}}.) To share? (\textbf{\ru{делиться}}.) To criticise? (\textbf{\ru{критиковать}}.) Criticism? (\textbf{\ru{критика}}.) To approve of? (\textbf{\ru{одобрять}}.) To condemn? (\textbf{\ru{осуждать}}.) To reproach? (\textbf{\ru{упрекать}}.) Logic? (\textbf{\ru{логика}}.) Logical? (\textbf{\ru{логичный}}.) A conclusion? (\textbf{\ru{вывод}}.) To reason? (\textbf{\ru{рассуждать}}.) Obvious? (\textbf{\ru{очевидный}}.)
\end{tcolorbox}
